\section{Conclusion}
\label{sec:conclusion}

We asked whether a quantum kernel used only during training helps a lightweight waste classifier beyond what a classical kernel in the same place provides.
We tested this on TrashNet, with RPS-25 as a low-data confirmation, using three students and five methods under one pipeline.
A projected quantum kernel was not distinguishable from logit KD in accuracy. It was ahead on the mean in \KdAheadFp{} of \NPairs{} pairs and behind on both TrashNet MobileNets, and no paired difference was significant.
At its pre-specified scale it beat its RBF twin in \GtwoFixedWins{} of \NPairs{} pairs.
Its seed-to-seed variance tended to be lower than the twin's and, less consistently, than that of hard-label training. This tendency is not significant at three seeds, depends on the metric and on one mismatched run, and does not hold on every pair.
Removing the circuit's entanglement hurt on TrashNet validation and helped on RPS-25, and fitting the teacher projection did not help.
All \NIntModels{} exported MobileNet students are fully integer INT8 models of about \MbMnvTwo{} and \MbMnvThree{}\,MB. They run in \LatMnvTwo{} and \LatMnvThree{}\,ms per inference on a Raspberry Pi~4B, and the choice of INT8 kernels on the device changes between \XnnDisMin{}\% and \XnnDisMax{}\% of a model's predictions.
The quantum kernel leaves nothing in the deployed model, and our evidence does not show an accuracy benefit over classical alternatives.

These results are exploratory.
One question remains open: whether the quantum map beats classical maps of equal width. Answering it would need a frozen protocol, the random-feature control we did not run, at least five seeds and fresh data. On the deployment side, energy and further devices remain to be measured.
